% Ad Atticum 15.26 --- headnote
% ad-atticum-15-26, 2 July 44 BC, in Arpinati

\textit{Cicero to Atticus from the Arpinum estate, 2 July 44 BC
--- Perseus dateline \textit{Scr. in Arpinati vi Non. Quint. a.
710 (44)}. Cicero is moving south down the Latin coast towards
Puteoli, the port from which he means to sail; he has paused at
Arpinum, the family estate. The letter is a five-section omnibus
of business: Quintus's family entanglement with Lepta and his son;
a rumour that Lucius Piso means to slip out of Italy on a forged
senatorial decree; the courier who reached him at Anagnia carrying
yet another letter from Brutus asking him to attend Brutus's games
in Rome; his refusal to do so; the games as a touchstone of
political nerve; and the various small matters of property he has
left in Atticus's hands --- Marcus Aelius's water-rights, the
half-share of the Tullianum that Cascellius has been pressing him
on, the eighth-share at \dag 380,000\dag\ sesterces.}

\textit{The political weight of the games is the heart of section
1. Brutus, in exile from the city in all but name, is required by
his praetorship to put them on; he must do so through delegates,
and his political survival turns on whether the Roman crowd
receives the festival warmly. Cicero, who has been refusing to set
foot in Rome since the arms of Antony began, will not attend ---
to do so would be neither necessary nor (he insists) honourable
--- but he is desperately invested in the outcome, and asks
Atticus for daily reports from the opening commission onward. The
Greek phrases are typical of the genre: \textit{pseudengraph\=oi}
(a forged decree), \textit{atop\=otaton} (the most absurd thing in
the world). The daggered clause in section 4 on the eighth share
is corrupt; the figures and the broker's name are clear, but the
syntax has not survived intact.}
